\subsection{Zariski rings}

PROPOSITION 6. Let A be a commutative Noetherian ring and m an ideal of A. The following properties are equivalent:

\begin{enumerate}
\def\labelenumi{(\alph{enumi})}
\item
  \(\mathfrak{m}\) is contained in the Jacobson radical of A.
\item
  Every finitely generated A-module is Hausdorff with the m-adic topology.
\item
  For every finitely generated A-module E, every submodule of E is closed with respect to the m-adic topology on E.
\item
  Every maximal ideal of A is closed with respect to the m-adic topology.
\end{enumerate}

Let us show that (a) implies (b). Suppose that m is contained in the Jacobson radical of A and let E be a finitely generated A-module. If \(x \in E\) and \(m \in m\) are such that \((1 - m)x = 0\) , then x = 0, for 1 - m is invertible in A. Then (no. 2, Proposition 5) E is Hausdorff with the m-adic topology.

Let us prove that (b) implies (c). Suppose (b) holds. Let E be a finitely generated A-module and F a submodule of E. Then E/F is Hausdorff with the m-adic topology, which is the quotient topology of the m-adic topology on E; then F is closed in E.

Clearly (c) implies (d). Let us show finally that (d) implies (a). It follows from (d) that, for every maximal ideal a of A, the A-module A/a is Hausdorff with the m-adic topology. This implies \(\mathfrak{m}(\mathrm{A}/\mathfrak{a}) \neq \mathrm{A}/\mathfrak{a}\) , unless the m-adic topology on A/a were the coarsest topology and A/a were reduced to 0, which is absurd since A/a is a field. The canonical image of m in A/a is therefore an ideal of A/a distinct from A/a and hence reduced to 0; then \(m \subset a\) , which proves that m is contained in the Jacobson radical of A.

DEFINITION 2. A topological ring A is called a Zariski ring if it is commutative and Noetherian and there exists a defining ideal m for the topology on A satisfying the equivalent conditions of Proposition 6.

A Zariski ring A is necessarily Hausdorff (Proposition 6) and every defining ideal of its topology is contained in the Jacobson radical of A.

\textbf{Examples of Zariski rings}\par

\begin{enumerate}
\def\labelenumi{(\arabic{enumi})}
\item
  Let A be a commutative Noetherian ring and m an ideal of A. If A is Hausdorff and complete with the m-adic topology, A is a Zariski ring with this topology, by virtue of § 2, no. 13, Lemma 3.
\item
  Every quotient ring A/b of a Zariski ring is a Zariski ring, for it is Noetherian and, if m is a defining ideal of A, \(\mathfrak{m}(\mathbf{A}/\mathfrak{b}) = (\mathfrak{m} + \mathfrak{b})/\mathfrak{b}\) is contained in the Jacobson radical of A/b (Algebra, Chapter VIII, § 6, no. 3, Proposition 7).
\item
  Let A be a Noetherian semi-local ring and r its Jacobson radical. Then A with the r-adic topology is a Zariski ring. This will always be the topology in question (unless otherwise stated) when we consider a Noetherian semi-local ring as a topological ring.
\end{enumerate}

PROPOSITION 7. Let A, A' be two commutative rings and h: A → A' a ring homomorphism. Suppose that A is Noetherian and that A' is a finitely generated A-module (with the structure defined by h). Let m be an ideal of A and let m' = mA'. Then:

\begin{enumerate}
\def\labelenumi{(\roman{enumi})}
\item
  For the \(\mathfrak{m}'\) -adic topology on \(\mathbf{A}'\) to be Hausdorff, it is necessary and sufficient that the elements of \(1 + h(\mathfrak{m})\) be not divisors of 0 in \(\mathbf{A}'\) .
\item
  If A with the m-adic topology is a Zariski ring, then A' with the m'-adic topology is a Zariski ring.
\item
  If \(h\) is injective (thus identifying \(\mathbf{A}\) with a subring of \(\mathbf{A}'\) ), the \(\mathfrak{m}'\) -adic topology on \(\mathbf{A}'\) induces on \(\mathbf{A}\) the \(\mathfrak{m}\) -adic topology.
\end{enumerate}

Recall that the m'-adic filtration on A' coincides with the m-adic filtration on the A-module A' (§ 2, no. 1, Example 3). Assertion (i) is thus a special case of Proposition 5 of no. 2 and assertion (iii) a special case of Theorem 2 of no. 2, Finally let us show (ii). Suppose that A is a Zariski ring with the m-adic topology and let E' be a finitely generated A'-module; it is also a finitely generated A-module and the m-adic and m'-adic filtrations on E' coincide; then E' is Hausdorff with the m'-adic topology. Finally the A-module A' is Noetherian and hence the ring A' is Noetherian, which completes the proof that A' is a Zariski ring.

